\chapter{DPF HMC Target Correctness and Structural-Model Suitability}
\label{ch:bf-dpf-hmc-target-suitability}

The preceding particle, transport, and learned-surrogate chapters built a ladder
of increasingly differentiable filtering constructions: raw particle-flow
proposals, proposal-corrected PF-PF, soft differentiable resampling,
deterministic OT equal-weighting, entropic optimal-transport relaxation with
Sinkhorn, custom-gradient LEDH-PFPF-OT differentiation for the declared finite
value path, retained-teacher learned transport acceleration, and broader
learned-map or dynamic-transport surrogates.
This chapter asks a narrower and more severe
question.  If one of those constructions is used inside Hamiltonian Monte Carlo
for a nonlinear DSGE or MacroFinance model, what target is the Markov chain
intended to leave invariant?

The answer cannot be inferred from the label ``differentiable particle
filter.''  HMC is a target-specific algorithm.  Its mathematical interpretation
depends on the scalar log target used for accept/reject decisions, the gradient
used by the numerical integrator, and the correction mechanism, if any
\citep{neal2011mcmc,betancourt2017conceptual}.  Particle MCMC adds a separate
target logic: an unbiased, nonnegative likelihood estimator can define a
corrected extended-space chain under pseudo-marginal assumptions, but that is
not the same object as a smooth surrogate likelihood differentiated by HMC
\citep{andrieu2009pseudo,andrieu2010particle}.  The DPF ladder must therefore
separate three questions:
\begin{enumerate}[label=(\roman*)]
  \item whether the scalar value supplied to HMC is the intended posterior
    value, an unbiased estimator inside a corrected construction, or a
    deliberately approximate value;
  \item whether the gradient supplied to the integrator is the derivative of
    that same scalar value;
  \item whether the evidence is strong enough for exploratory, research, or
    broader external claims.
\end{enumerate}

\section{Objects in the HMC target contract}
\label{sec:bf-dpf-hmc-contract}

Let $\theta\in\Theta$ be the structural parameter and let
$u\in\mathbb{R}^{d_\theta}$ be the unconstrained parameter used by the sampler,
with transform $\theta=T(u)$.  Let
\[
  \ell(u)
  =
  \log p(T(u))+\log L(y_{1:T}\mid T(u))+\log\left|\det DT(u)\right|
\]
denote the scalar log target in unconstrained coordinates when the filtering
likelihood is exact.  This expression assumes that the transform is
differentiable on the sampler-valid region and that $DT(u)$ is the square
Jacobian relevant to the change of variables, or the appropriate constrained-
to-unconstrained volume term if the parameterization is more general.  In a DPF
construction the likelihood term may be replaced by a particle, relaxed, or
learned object.  The HMC contract is therefore stated for the actual scalar
value $\ell_\star(u)$ used by the sampler, not for the method family name.

\begin{longtable}{@{}p{0.18\textwidth} p{0.38\textwidth} p{0.36\textwidth}@{}}
\caption{Objects that must be named before a DPF construction can be called an
HMC target.}
\label{tab:bf-dpf-hmc-object-inventory}\\
\toprule
Object & Mathematical role & Reviewer audit question \\
\midrule
\endfirsthead
\toprule
Object & Mathematical role & Reviewer audit question \\
\midrule
\endhead
Parameter $u$ and transform $T$ &
Coordinates integrated by HMC and the map back to structural parameters. &
Are support constraints, Jacobian terms, and invalid parameter regions included
in the scalar target? \\
Scalar log target $\ell_\star(u)$ &
The value whose density $\pi_\star(u)\propto\exp\{\ell_\star(u)\}$ is claimed
as the invariant target. &
Is $\ell_\star$ exact, pseudo-marginal extended-space, relaxed, or learned? \\
Potential $U_\star(u)$ &
$U_\star(u)=-\ell_\star(u)$, up to an irrelevant additive constant. &
Is the sign convention consistent between value, gradient, and accept/reject? \\
Momentum $p$ and mass $M$ &
Auxiliary Gaussian momentum $p\sim\mathcal{N}(0,M)$ used to define Hamiltonian
dynamics. &
Is $M$ fixed or adapted under a declared warmup policy, and is it positive
definite? \\
Hamiltonian $H_\star(u,p)$ &
$H_\star(u,p)=U_\star(u)+\frac{1}{2}p^\top M^{-1}p$ in ordinary separable HMC. &
Does the accept/reject step evaluate this same $H_\star$? \\
Integrator $\Phi_{\epsilon,L}$ &
Numerical map used to propose $(u',p')$ after $L$ steps of size $\epsilon$. &
Is the proposal reversible and volume preserving for the stated correction, or
is a different Metropolis ratio supplied? \\
Gradient $g(u)$ &
Vector supplied to the integrator in place of $\nabla_u\ell_\star(u)$. &
Is $g(u)=\nabla_u\ell_\star(u)$ for the same scalar value, or is the mismatch
explicitly corrected and diagnosed? \\
Accept/reject value &
Scalar value used in the Metropolis correction. &
Does the correction use the intended target, the relaxed target, or a learned
surrogate? \\
Target-status category &
Label attached to the resulting chain. &
Does the label distinguish exact, pseudo-marginal, relaxed, learned, and
diagnostic-only use? \\
\bottomrule
\end{longtable}

Assume that $\ell_\star:\mathbb{R}^{d_\theta}\to\mathbb{R}$ is differentiable
on the interior of the sampler's valid region, and that
$g:\mathbb{R}^{d_\theta}\to\mathbb{R}^{d_\theta}$ is the vector field supplied
to the integrator in the same coordinates.  The mass matrix $M$ in the ordinary
separable Hamiltonian is a symmetric positive-definite
$d_\theta\times d_\theta$ matrix, so $p^\top M^{-1}p$ is well defined.  The
minimum value-gradient consistency condition is
\begin{equation}
\label{eq:bf-dpf-hmc-contract}
  g(u) = \nabla_u \ell_\star(u).
\end{equation}
Equation~\eqref{eq:bf-dpf-hmc-contract} is not a claim that
$\ell_\star$ is the exact nonlinear filtering likelihood.  It is a claim that
the gradient used by the integrator and the value used by the target
interpretation are the same mathematical object.

If the value is $\ell_\star$ but the integrator uses a vector field
$g\ne\nabla\ell_\star$, the usual leapfrog derivation no longer describes the
Hamiltonian flow of $H_\star$.  A Metropolis correction may still be valid for
an arbitrary proposal only after reversibility, volume preservation, and the
correct acceptance ratio have been established.  Without that separate
proposal-level proof, a smooth chain with mismatched value and gradient is not
evidence of correct HMC.  If both value and gradient are changed to a relaxed
or learned scalar, then the chain may be internally coherent, but the target has
changed to $\pi_\star$.

\section{Target-status taxonomy}
\label{sec:bf-dpf-hmc-taxonomy}

The DPF chapters use the following target-status categories.  The categories
are deliberately conservative because a reviewer should be able to tell exactly
what posterior, extended posterior, or surrogate posterior is being discussed.

\begin{description}[leftmargin=0.18\textwidth, style=nextline]
  \item[Exact-likelihood HMC]
  HMC evaluates a deterministic scalar log posterior and its exact or audited
  derivative.  Numerical integration error is corrected by the standard
  Metropolis step, provided the implemented proposal satisfies the usual
  HMC reversibility and volume-preservation assumptions.

  \item[Pseudo-marginal MCMC]
  A random nonnegative likelihood estimator is included as part of an
  extended-space target.  The claim is not that the estimator is a smooth
  likelihood surrogate.  The claim is that marginal correctness follows from
  the pseudo-marginal construction under its assumptions
  \citep{andrieu2009pseudo,andrieu2010particle}.

  \item[Noisy-gradient method]
  A stochastic or approximate gradient is used to move the chain.  This is not
  exact HMC unless a valid correction theory is stated.  It may be useful as an
  optimization, adaptation, or diagnostic device, but that usefulness does not
  by itself identify the invariant posterior.

  \item[Delayed-acceptance or surrogate-corrected MCMC]
  A surrogate or particle object proposes, screens, or accelerates moves, but a
  later correction evaluates the intended target.  The target claim belongs to
  the corrected chain, not to the surrogate alone.

  \item[Relaxed-target HMC]
  The scalar log target deliberately replaces a discontinuous or discrete
  filtering operation by a smooth relaxed operation, such as soft resampling or
  entropic OT.  HMC can be coherent for this scalar if
  Equation~\eqref{eq:bf-dpf-hmc-contract} holds, but the posterior is the
  relaxed posterior.

  \item[Learned-surrogate HMC]
  A neural or amortized map defines part of the scalar value.  The chain targets
  the learned scalar only if value and gradient are consistent.  It targets the
  original or teacher posterior only after a separate correction or error-bound
  argument strong enough for that claim.
\end{description}

\section{Rung-by-rung target analysis}
\label{sec:bf-dpf-hmc-rung-analysis}

\subsection{EDH and LEDH flow proposals}

At the raw EDH or LEDH rung, the value object is a flow-based approximate
filtering quantity.  The flow is derived from Gaussian closure or local
linearization assumptions and is integrated numerically.  If HMC differentiates
the same scalar flow object that it evaluates, then the construction can satisfy
value-gradient consistency for that scalar.  Its target status is nevertheless
limited:
\begin{itemize}
  \item exact-target HMC only in recovery cases where the flow construction
    equals the exact filtering likelihood, such as the linear-Gaussian
    benchmark;
  \item approximate-target or value-side benchmark HMC outside those recovery
    cases;
  \item diagnostic-only use if the implementation differentiates a path that
    differs from the scalar value used for accept/reject.
\end{itemize}
Smoothness of the flow ODE, by itself, does not upgrade a closure-based
filtering approximation into an exact nonlinear posterior target.

\subsection{PF-PF proposal correction}

At the PF-PF rung, the flow is used as a proposal and the scalar value includes
the proposal-to-target correction through importance weights and
change-of-variables terms.  This is the first rung in the BayesFilter DPF ladder
with an explicit proposal-correction interpretation.  That statement is not a
production or validation claim.  It means only that the value-side target story
is more defensible than raw flow transport because the proposal density,
Jacobian or log determinant, and particle weight enter the same scalar object.
The improvement is proposal-accounting clarity for the stated one-step
conditional construction in Chapter~\ref{ch:bf-dpf-pfpf}; it is not evidence, by
itself, of multi-step posterior fidelity, finite-particle stability, or HMC
correctness.

For HMC, the required scalar is the corrected particle objective actually used
by the sampler.  The required gradient is the derivative of that same corrected
objective, including every differentiable term in the proposal density,
transition density, measurement density, flow map, and log determinant.  The
remaining evidence burden is substantial:
\begin{itemize}
  \item finite-particle variability must be measured rather than hidden by a
    deterministic seed;
  \item log-determinant and Jacobian paths must be audited against the same
    value path;
  \item flow discretization error must be separated from Monte Carlo error;
  \item compiled repeated evaluation must remain finite and reproducible on the
    structural model class being claimed;
  \item if randomness enters the HMC target, the construction must state
    whether it is pseudo-marginal extended-space MCMC, a fixed-randomness
    deterministic approximation, or a noisy-gradient diagnostic.
\end{itemize}
Thus PF-PF is a plausible first research rung for DPF-HMC experiments, not a
full HMC-validity theorem and not by itself a stronger external claim.

\subsection{Soft differentiable resampling}

Soft resampling replaces categorical ancestor selection by a differentiable
mixture or shrinkage construction \citep{zhumurphyjonschkowski2020,jonschkowski2018differentiable,karkus2018particle}.  The scalar value is therefore a
soft-resampling filtering value.  If HMC differentiates exactly that scalar, the
resulting chain can be internally coherent for the soft target.  It should not
be described as sampling the categorical-resampling posterior unless a separate
correction is supplied.

The core mathematical issue is nonlinear test-function bias.  Even if a
soft-resampling map preserves a weighted mean under special conditions, it does
not generally preserve expectations of nonlinear functions, path genealogies,
or the likelihood contribution of the original categorical particle filter.
The HMC label must therefore be relaxed-target HMC, surrogate-target HMC, or
diagnostic use, depending on the scalar value actually corrected by the
sampler.

\subsection{OT and entropic OT resampling}

At the OT/EOT rung, resampling is written as a transport problem \citep{reich2013nonparametric,cuturi2013sinkhorn,peyre2019computational,corenflos2021differentiable}.  Entropic
regularization and a finite Sinkhorn solve make the construction differentiable
and numerically structured, but they also define a different map from exact
categorical resampling.  The scalar value must specify the cost, entropy
convention, regularization parameter, stopping tolerance, and whether the
unrolled solver path or an implicit derivative is used.

For HMC, the target status is relaxed-target HMC when the accept/reject value
and gradient both use the same EOT/Sinkhorn scalar.  Finite solver residuals,
underflow handling, stabilization, and $\varepsilon$ sensitivity are target
diagnostics, not mere implementation details.  A finite Sinkhorn layer does not
become an exact categorical resampling layer because its gradient exists.

\subsection{Learned or amortized OT}

Learned OT introduces a second approximation layer: a neural or amortized map is
trained to imitate, accelerate, or residual-correct a transport teacher
\citep{amos2017icnn,makkuva2020otmapping,uscidda2023mongegap,gracykchen2024geonet}.  The
scalar value is then the learned scalar unless a correction evaluates the
teacher or original target.  The gradient path is the derivative through the
learned map and any explicitly retained teacher terms.

The target-status label must therefore be learned-surrogate HMC unless the
implementation supplies a valid correction or an error analysis strong enough
to promote the claim.  Runtime reduction, smoothness, and low teacher residual
on training clouds are not target-correctness evidence.  The missing validation
includes out-of-distribution particle-cloud tests, seed sensitivity, posterior
comparison against the EOT teacher, and stress tests on structural parameter
regions that are invalid, nearly invalid, or geometrically stiff.

\section{Structured target maps}
\label{sec:bf-dpf-hmc-summary}

\begin{longtable}{@{}p{0.11\textwidth} p{0.20\textwidth} p{0.20\textwidth} p{0.16\textwidth} p{0.25\textwidth}@{}}
\caption{Rung-by-rung HMC target map.}
\label{tab:bf-dpf-hmc-rungs}\\
\toprule
Rung & Scalar value & Gradient path & Target status & Missing validation before promotion \\
\midrule
\endfirsthead
\toprule
Rung & Scalar value & Gradient path & Target status & Missing validation before promotion \\
\midrule
\endhead
EDH/LEDH &
Gaussian-closure or local-flow filtering scalar. &
Derivative of the same flow scalar, if implemented consistently. &
Exact only in recovery cases; otherwise approximate benchmark. &
Linear-Gaussian recovery, nonlinear likelihood-error studies, flow
discretization checks, invalid-region behavior. \\
PF-PF &
Proposal-corrected particle scalar with weights and flow Jacobian/log-det terms. &
Derivative of the same corrected scalar, including correction terms. &
First rung with explicit proposal-correction interpretation; research
candidate only. &
Finite-particle variance, Jacobian/log-det audit, value-gradient parity,
compiled repeated evaluation, randomness policy. \\
Soft resampling &
Soft or shrinkage resampling scalar. &
Derivative through the same soft map. &
Relaxed or surrogate target. &
Posterior sensitivity against categorical or trusted references, nonlinear
test-function bias, seed and temperature sensitivity. \\
OT/EOT &
Transport-relaxed scalar with cost, entropy, $\varepsilon$, and solver budget. &
Derivative through the same unrolled or implicit solver path. &
Relaxed target. &
$\varepsilon$ and Sinkhorn-budget sensitivity, residual tolerances,
stabilization audit, teacher posterior comparison. \\
Learned OT &
Learned approximation to an OT/EOT or residual-corrected scalar. &
Derivative through the learned map and retained teacher terms. &
Learned surrogate unless corrected. &
Teacher residual to posterior-error analysis, out-of-distribution cloud tests,
structural stress tests, correction evidence. \\
\bottomrule
\end{longtable}

Table~\ref{tab:bf-dpf-hmc-rungs} is a target map, not a performance table.  It
does not say that the lower rungs are useless or that the higher rungs are
invalid.  It says that every rung must name the scalar target before an HMC
claim can be interpreted.

\section{Pseudo-marginal and surrogate distinctions}
\label{sec:bf-dpf-hmc-pm-surrogate}

The pseudo-marginal result and the DPF surrogate story are often confused
because both involve particle approximations.  They are different claims.
Suppose $\widehat{L}(y\mid\theta,\omega)\ge 0$ is an unbiased estimator of
$L(y\mid\theta)$ under auxiliary randomness $\omega$:
\[
  \mathbb{E}_{\omega\mid\theta}\left[\widehat{L}(y\mid\theta,\omega)\right]
  =
  L(y\mid\theta).
\]
Pseudo-marginal MCMC targets an extended density involving $\omega$ and can
recover the exact marginal posterior under its assumptions.  The validity
argument is an extended-target argument.  It is not an argument that
$\log\widehat{L}$ is a smooth deterministic likelihood whose pathwise gradient
is the gradient of $\log L$.

By contrast, a relaxed DPF target usually chooses a deterministic or
fixed-randomness scalar $\ell_\star(\theta)$ and differentiates it.  HMC may be
valid for $\ell_\star$ if the value-gradient contract and proposal correction
hold.  That does not imply validity for the original posterior
$\ell(\theta)$.  A delayed-acceptance or surrogate-corrected scheme can bridge
the gap only if the final correction evaluates the intended target with the
right acceptance probability.

\section{Why nonlinear DSGE models sharpen the issue}
\label{sec:bf-dpf-hmc-dsge}

Nonlinear DSGE models make target drift especially costly because the filtering
law is tied to structural economics rather than only to a generic latent-state
transition.  The DPF chapter does not need to solve the economics, but it must
not hide target changes inside a numerical method.  The main DSGE stress points
are:
\begin{enumerate}[label=(\roman*)]
  \item \textbf{Determinacy and invalid regions}: the sampler will propose
    parameters outside or near determinacy regions.  The target must return a
    declared finite penalty or $-\infty$ policy, not an uncaught solver failure.
  \item \textbf{Pruning and approximation convention}: first-order, higher-order,
    pruned, and simulation-based approximations can define different latent
    laws.  A DPF likelihood is only meaningful after the structural law is
    fixed.
  \item \textbf{Structural timing}: observables, controls, states, and shocks
    must enter the transition and measurement equations with the intended
    timing.  A smooth filter for the wrong timing is still the wrong target.
  \item \textbf{Latent-state semantics}: deterministic-completion coordinates
    and stochastic state coordinates are not interchangeable.  Particle clouds
    must represent the declared latent law.
  \item \textbf{Local Jacobian fragility}: EDH, LEDH, and proposal-corrected
    flow methods inherit local-linearization and derivative fragility near
    boundaries or stiff regions.
\end{enumerate}

For DSGE work, the correct conclusion is conservative.  A DPF rung may be a
useful research target after it states its scalar value and diagnostics.  It is
not bank-facing evidence until the structural law, invalid-region policy,
value-gradient parity, and posterior sensitivity have been tested on the named
model class.

\section{Why MacroFinance models sharpen the issue}
\label{sec:bf-dpf-hmc-macrofinance}

MacroFinance latent-factor systems create a different target-risk profile.
Their long panels can amplify small likelihood changes, and their latent
factors often interact with volatility, measurement noise, and missing-data
patterns.  A relaxed or learned resampling map that looks harmless on a short
toy filter can move posterior mass materially in a long panel.

The minimum MacroFinance limitations to record are:
\begin{itemize}
  \item latent dimension and particle degeneracy pressure;
  \item long observation spans that compound small per-period target changes;
  \item missing, ragged-edge, or mixed-frequency observation patterns;
  \item volatility and measurement-noise directions with sharp curvature;
  \item compiled repeatability across seeds, batches, and hardware settings;
  \item model-risk governance, including reproducible artifacts and claim
    boundaries.
\end{itemize}
These limitations do not rule out DPF research.  They rule out treating a
smooth differentiable filtering path as banking evidence before the target and
validation ladder have been passed.

\section{Relation to surrogate-HMC acceleration}
\label{sec:bf-dpf-hmc-hnn-surrogates}

Surrogate-HMC and Hamiltonian-neural-network ideas address a different point in
the inference stack \citep{greydanus2019hamiltonian,hoffman2019neutra,parno2018transport}.  They attempt to improve
movement through an already chosen target, or to use a surrogate proposal that
is later corrected against that target.  DPF methods may instead alter the
filtering likelihood itself through flow approximation, relaxed resampling, or
learned transport.  The distinction is operational:
\begin{itemize}
  \item a corrected surrogate-HMC proposal can preserve the intended target if
    the correction evaluates that target and the proposal assumptions hold;
  \item a relaxed DPF likelihood defines a new target unless a correction
    returns to the original likelihood;
  \item a learned DPF transport map defines a learned target unless the sampler
    corrects against the teacher or original target.
\end{itemize}

\begin{longtable}{@{}p{0.18\textwidth} p{0.25\textwidth} p{0.23\textwidth} p{0.24\textwidth}@{}}
\caption{Method-family distinctions relevant to DPF-HMC claims.}
\label{tab:bf-dpf-hmc-family-comparison}\\
\toprule
Method family & What is approximated? & Target implication & Required wording \\
\midrule
\endfirsthead
\toprule
Method family & What is approximated? & Target implication & Required wording \\
\midrule
\endhead
Classical exact HMC &
Hamiltonian trajectory by numerical integration. &
Target preserved after Metropolis correction when value, gradient, and proposal
conditions hold. &
Exact for the named scalar target, not exact because the trajectory is exact. \\
Pseudo-marginal MCMC &
Likelihood value by an unbiased nonnegative estimator on extended space. &
Marginal posterior can be exact under pseudo-marginal assumptions. &
Corrected extended-space chain, not smooth surrogate-gradient HMC. \\
Noisy-gradient method &
Gradient or force field used for movement. &
No exact posterior claim without correction theory. &
Diagnostic or approximate unless corrected. \\
Delayed-acceptance surrogate &
Early-stage value or proposal geometry. &
Can preserve intended target if final correction evaluates intended target. &
Surrogate-assisted, not target-defining, when corrected. \\
Relaxed DPF &
Filtering likelihood construction itself. &
Posterior changes to relaxed posterior. &
Relaxed-target HMC. \\
Learned DPF &
Transport, residual, or filtering map by a learned function. &
Posterior changes to learned posterior unless corrected. &
Learned-surrogate HMC pending correction or validation. \\
\bottomrule
\end{longtable}

\section{Evidence ladder for stronger external language}
\label{sec:bf-dpf-hmc-banking-gates}

The chapter's recommendation is evidence-gated.  A method can be mathematically
interesting, credible as a research prototype, and still not yet justify
stronger external claims.  The following ladder constrains the wording used in
this chapter and in later experiments.

\begin{longtable}{@{}p{0.18\textwidth} p{0.24\textwidth} p{0.27\textwidth} p{0.21\textwidth}@{}}
\caption{Evidence ladder for DPF-HMC claim language.}
\label{tab:bf-dpf-hmc-banking-evidence}\\
\toprule
Claim level & Allowed claim & Required evidence & Disallowed shortcut \\
\midrule
\endfirsthead
\toprule
Claim level & Allowed claim & Required evidence & Disallowed shortcut \\
\midrule
\endhead
Exploratory use &
The rung is worth studying on controlled examples. &
Named scalar value, value-gradient check, finite smoke tests, clear target
status. &
Calling the method suitable because it is differentiable. \\
Credible research use &
The rung has a defensible target interpretation for a named model class. &
Source-grounded derivation, repeated fixed-seed and varied-seed tests,
posterior comparisons, failure-region diagnostics. &
Treating a toy result as evidence for DSGE or MacroFinance deployment. \\
Broader external claim &
The rung supports a restricted external research conclusion with stated
limitations. &
Reproducible demonstrations, model-specific stress tests, challenger
comparisons, and documented limitations. &
Using runtime improvement as target-correctness evidence. \\
Operational claim &
The rung is proposed for controlled use in a defined workflow. &
Independent review, change control, monitoring, rollback, data and model
stability across releases. &
Inferring operational readiness from chapter-level mathematics. \\
\bottomrule
\end{longtable}

No rung in this chapter is promoted to an operational claim.  The strongest
current recommendation is a bounded experimental hypothesis: prioritize
proposal-corrected PF-PF first because it is the first rung with an explicit
proposal-correction interpretation, then compare relaxed and learned rungs as
deliberately changed targets.  The hypothesis remains conditional on the finite-
particle, Jacobian, randomness, compilation, and posterior-comparison checks
stated above.

\section{Recommendation}
\label{sec:bf-dpf-hmc-recommendation}

The disciplined recommendation is:
\begin{enumerate}[label=(\roman*)]
  \item use EDH and LEDH as recovery and approximation benchmarks, not as
generic exact HMC targets;
  \item treat proposal-corrected PF-PF as the first experimental hypothesis with
    an explicit proposal-correction interpretation and a defensible value-side
    target story, subject to finite-particle, Jacobian, randomness,
    compilation, and posterior-comparison checks;
  \item label soft resampling and OT/EOT as relaxed-target HMC unless a
    correction restores the original categorical or trusted target;
  \item label learned OT as learned-surrogate HMC unless a correction or
    posterior-error argument justifies stronger wording;
  \item require model-class evidence before any stronger external claim.
\end{enumerate}

For the current repair lane, the serious experimental route is therefore not a
no-resampling shortcut.  The route to test is Li--Coates Algorithm~1 UKF LEDH
with PF-PF correction and a declared OT, Sinkhorn, or annealed-transport
resampling relaxation.  Carrying the Algorithm~1 covariance state through that
transport is an implementation contract for the selected route, because the
covariance state belongs to Algorithm~1 while the transport map defines the
relaxed resampling layer.  A no-resampling graph or fixed-randomness run can
still be useful as a kernel, shape, or gradient-plumbing diagnostic, but it is
not by itself evidence that the DPF route is filter-ready or HMC-ready.

\section{What this chapter does and does not claim}
\label{sec:bf-dpf-hmc-boundary}

This chapter claims that DPF-HMC analysis must be organized around named scalar
targets and value-gradient consistency.  It does not claim:
\begin{itemize}
  \item that PF-PF has already passed finite-particle, Jacobian, compiled-path,
    or posterior-comparison checks,
  \item that relaxed or learned resampling targets are unusable,
  \item that pseudo-marginal validity transfers to deterministic relaxed
    gradients,
  \item that nonlinear DSGE or MacroFinance DPF-HMC has already been fully
    validated.
\end{itemize}
Its role is narrower: it clarifies which scalar target is being differentiated,
what correction mechanism is being invoked if any, and what mathematical
interpretation the resulting chain can support.

\FloatBarrier
